% Rapport FINAL à rendre (~5 pages) : résumé du rapport complet (rapport.tex).
% Le détail, les justifications et les renvois vers le code sont dans rapport.tex ;
% ici, seulement l'essentiel, avec renvoi au dépôt pour le reste.
%
% Compiler depuis le dossier report/ :   latexmk -pdf rapport_final.tex
%
% Les figures sont partagées avec le rapport complet (figures/). Sections : sections_final/.

\documentclass[11pt,a4paper]{article}

\newif\ifsuivi
\suivifalse  % pas d'encadrés de travail dans le rapport rendu

% Préambule commun au rapport complet (rapport.tex) et au rapport final
% (rapport_final.tex). Le fichier qui l'inclut doit définir \ifsuivi avant.

\usepackage[utf8]{inputenc}
\usepackage[T1]{fontenc}
\usepackage[french]{babel}
\usepackage{lmodern}
\usepackage[hmargin=1.8cm, top=1.2cm, bottom=1cm, includehead, includefoot, headsep=8pt, footskip=16pt]{geometry}
\usepackage{amsmath}
\usepackage{graphicx}
\usepackage{booktabs}
\usepackage[section]{placeins} % les tableaux restent dans leur section
\usepackage{tabularx}
\usepackage{array}
\usepackage{enumitem}
\usepackage[dvipsnames]{xcolor}
\usepackage{siunitx}
\usepackage{fancyhdr}
\usepackage[font=small, labelfont=bf]{caption}
\usepackage[hidelinks]{hyperref}

\sisetup{locale = FR}
\graphicspath{{figures/}}
\setlist{nosep, leftmargin=*}
\newcolumntype{Y}{>{\raggedright\arraybackslash}X}
\renewcommand{\arraystretch}{1.15}

% Encadré de travail, visible seulement en version de suivi. Contenu : une liste
% \item ... de ce qui reste à faire dans la section.
\newcommand{\afaire}[1]{\ifsuivi\par\medskip\noindent
  \fcolorbox{BrickRed}{BrickRed!6}{\parbox{\dimexpr\linewidth-2\fboxsep-2\fboxrule}{\small
    \textcolor{BrickRed}{\textbf{À faire}}
    \begin{itemize}#1\end{itemize}}}\par\medskip\fi}
% Renvoi vers le code qui produit les chiffres d'une partie (version de suivi seulement).
\newcommand{\outrouver}[1]{\ifsuivi\par\smallskip\noindent
  {\small\sloppy\color{NavyBlue}\textbf{Où trouver :} #1\par}\smallskip\fi}
% Nom de colonne / de fichier du dépôt.
\newcommand{\code}[1]{\texttt{\detokenize{#1}}}
% Autorise une coupure de ligne après / et - dans les noms longs (sans tiret ajouté).
% Attention : les espaces sont ignorés ; pour un nom avec espace, utiliser \texttt.
\renewcommand{\code}[1]{\texttt{\codecut#1\codeend}}
\def\codeend{\codeend}
\def\codecut#1{\ifx#1\codeend\else\ifx#1/\slash\else\ifx#1-\char`\-\allowbreak\else\detokenize{#1}\fi\fi\expandafter\codecut\fi}


\pagestyle{fancy}
\fancyhf{}
\lhead{\small Prédiction de qualité de soudures}
\rhead{\small Équipe \numequipe}
\cfoot{\thepage}

\newcommand{\numequipe}{??} % TODO numéro d'équipe

\begin{document}

% ---------------------------------------------------------------- page de titre
% Le sujet demande : noms, numéro d'équipe, date, mails et titre en première page.
\begin{center}
  {\LARGE\bfseries Prédiction de qualité de soudures}\\[6pt]
  {\large Projet Machine Learning --- CentraleSupélec, 3A}\\[14pt]
  \begin{tabular}{c@{\hspace{2.5em}}c}
    Pablo Bosch & Myriam Boulaares \\
    \footnotesize\texttt{?} & \footnotesize\texttt{myriam.boulaares@student-cs.fr} \\[6pt]
    Salomé Fonvielle & Margaux Langlois \\
    \footnotesize\texttt{salome.fonvielle@gmail.com} & \footnotesize\texttt{?}
  \end{tabular}\par\vspace{8pt}
  Équipe n\textsuperscript{o}~\numequipe \quad --- \quad \today
\end{center}
\thispagestyle{fancy}

\section{Introduction}

% À rédiger à partir du rapport complet (rapport.tex).

\section{Données et pré-traitement}

% À rédiger à partir du rapport complet (rapport.tex).
% Résumer sections 2.1 et 2.2 du rapport complet : base, constats clés (valeurs manquantes MAR, une ligne = un essai, Charpy), choix de pré-traitement ; 1 ou 2 figures (ACP, corrélations).

\section{Méthodes et protocole de validation}

% À rédiger à partir du rapport complet (rapport.tex).

\section{Résultats et comparaison}

% À rédiger à partir du rapport complet (rapport.tex).

\section{Conclusion et recommandations}

% À rédiger à partir du rapport complet (rapport.tex).


\end{document}
